xelatex emr_komissiya.tex && xelatex emr_komissiya.tex
% !TEX program = xelatex
\documentclass[12pt,a4paper]{article}

\usepackage{fontspec}
\usepackage{polyglossia}
\setmainlanguage{azerbaijani}
\setmainfont{Georgia}
\newfontfamily\azfont{Georgia}

\usepackage[a4paper,left=30mm,right=15mm,top=20mm,bottom=20mm]{geometry}
\usepackage{setspace}
\usepackage{enumitem}
\usepackage{array}
\usepackage{tabularx}
\usepackage{ragged2e}
\usepackage{titlesec}
\usepackage{fancyhdr}

\setstretch{1.15}
\sloppy
\setlength{\emergencystretch}{3em}
\hyphenpenalty=500
\tolerance=2000
\setlength{\parindent}{0pt}
\setlength{\parskip}{0pt}

\pagestyle{fancy}
\fancyhf{}
\renewcommand{\headrulewidth}{0pt}
\fancyfoot[C]{\small\thepage}

\newcommand{\xett}{\rule{45mm}{0.4pt}}
\newcommand{\adxetti}{\rule{52mm}{0.4pt}}
\newcommand{\imzaxetti}{\rule{40mm}{0.4pt}}

\begin{document}

\begin{center}
  {\large\bfseries AZƏRBAYCAN RESPUBLİKASININ}\\[2pt]
  {\large\bfseries TƏHSİL İNSTİTUTU}\\[10pt]
  \rule{\textwidth}{0.8pt}
\end{center}

\vspace{6pt}

\begin{center}
  {\Large\bfseries Ə\,M\,R}
\end{center}

\vspace{10pt}

\begin{flushleft}
\begin{tabularx}{\textwidth}{@{}Xr@{}}
  \xett\ \ «\ \ \ \ \ \ » \ \ 2026-cı il & №\ \xett \\
  Bakı şəhəri & \\
\end{tabularx}
\end{flushleft}

\vspace{14pt}

\begin{center}
  {\bfseries Süni intellektlə idarəetmənin təşkili üzrə komissiyanın\\
  yaradılması haqqında}
\end{center}

\vspace{12pt}

\justifying

Azərbaycan Respublikasının Təhsil İnstitutunun fəaliyyətinin süni intellekt
texnologiyaları vasitəsilə idarə olunmasının təşkili, struktur bölmələrinin
məlumat resurslarının vahid idarəetmə sisteminə inteqrasiyası, idarəetmə
qərarlarının qəbulunda analitik imkanların genişləndirilməsi və bu sahədə
işlərin mərhələli şəkildə həyata keçirilməsi məqsədilə, Azərbaycan
Respublikasının Təhsil İnstitutunun Nizamnaməsinə uyğun olaraq,

\vspace{8pt}
\begin{center}
  {\bfseries Ə\,M\,R\ \ E\,D\,İ\,R\,Ə\,M:}
\end{center}
\vspace{8pt}

\begin{enumerate}[leftmargin=18pt,itemsep=6pt,topsep=2pt]

\item Azərbaycan Respublikasının Təhsil İnstitutunun fəaliyyətinin süni
intellektlə idarə olunmasının təşkili üzrə Komissiya (bundan sonra ---
Komissiya) yaradılsın və onun tərkibi \textbf{Əlavə № 1}-ə uyğun təsdiq
edilsin.

\item Müəyyən edilsin ki, Komissiya aşağıdakı vəzifələri həyata keçirir:

  \begin{enumerate}[label=\arabic{enumi}.\arabic*.,leftmargin=28pt,itemsep=3pt,topsep=3pt]
    \item institutun struktur bölmələrinin məlumat resurslarının
      (məlumat bazalarının, hesabat formalarının, statistik göstəricilərin)
      inventarizasiyasını aparmaq və onların vahid reyestrini tərtib etmək;
    \item süni intellekt əsaslı idarəetmə sisteminin texniki tələblərini,
      arxitekturasını və tətbiq mərhələlərini müəyyən etmək;
    \item məlumatların toplanması, emalı, saxlanması və qorunması qaydalarını,
      o cümlədən fərdi məlumatların mühafizəsi üzrə tələbləri hazırlamaq;
    \item pilot layihənin həyata keçirilməsi üzrə təkliflər paketini
      hazırlayıb institut rəhbərliyinə təqdim etmək;
    \item tətbiqin nəticələrinin qiymətləndirilməsi metodikasını və
      nəzarət göstəricilərini hazırlamaq.
  \end{enumerate}

\item İnstitutun mərkəz rəhbərləri, maliyyə və buxalteriya üzrə məsul
şəxsləri, habelə Funksional şöbələrin müdiri Komissiyanın fəaliyyəti üçün
zəruri olan məlumatların vaxtında, tam və düzgün təqdim edilməsini təmin
etsinlər.

\item Komissiyanın iclasları ayda azı bir dəfə keçirilsin; iclasların
protokolları və hesabat materialları Elmi katibliyə təqdim edilsin.

\item Komissiyanın fəaliyyətinin maddi-texniki, informasiya və kadr
təminatı Funksional şöbələrin müdiri tərəfindən həyata keçirilsin.

\item Komissiyaya bu əmrin icrası üçün zəruri olan əlavə məlumatları
institutun bütün struktur bölmələrindən almaq hüququ verilsin.

\item Bu əmr imzalandığı gündən qüvvəyə minir.

\item Əmrin icrası institutun direktor müavini
\textbf{Talıbov Tariyel İsmayıl oğluna} həvalə edilsin.

\item Əmrin icrasına nəzarəti öz üzərimdə saxlayıram.

\end{enumerate}

\vspace{22pt}

\begin{flushleft}
\begin{tabular}{@{}p{60mm}p{45mm}p{62mm}@{}}
  Direktor əvəzi & \imzaxetti & Əliyev Elnur Qəzənfər oğlu \\
\end{tabular}
\end{flushleft}

\vspace{16pt}

\begin{flushleft}
  \textbf{Viza:}\\[6pt]
  \begin{tabular}{@{}p{60mm}p{45mm}p{62mm}@{}}
    Elmi katib & \imzaxetti & Həsənova Aygün Vahid qızı \\
  \end{tabular}
\end{flushleft}

\newpage

\begin{center}
  {\bfseries Əlavə № 1}\\[4pt]
  Azərbaycan Respublikasının Təhsil İnstitutunun direktor əvəzinin\\
  «\ \ \ \ \ \ » \ \ 2026-cı il tarixli №\ \xett\ Əmri ilə\\
  \textbf{TƏSDİQ EDİLMİŞDİR}
\end{center}

\vspace{16pt}

\begin{center}
  {\bfseries TƏHSİL İNSTİTUTUNUN FƏALİYYƏTİNİN SÜNİ İNTELLEKTLE\\
  İDARƏ OLUNMASININ TƏŞKİLİ ÜZRƏ KOMİSSİYANIN\\
  T\,Ə\,R\,K\,İ\,B\,İ}
\end{center}

\vspace{14pt}

\begin{flushleft}
\textbf{Sədr:}\\[3pt]
\begin{tabular}{@{}p{68mm}p{96mm}@{}}
  Talıbov Tariyel İsmayıl oğlu & --- institutun direktor müavini \\
\end{tabular}

\vspace{10pt}
\textbf{Katib:}\\[3pt]
\begin{tabular}{@{}p{68mm}p{96mm}@{}}
  Həsənova Aygün Vahid qızı & --- elmi katib \\
\end{tabular}

\vspace{14pt}
\textbf{Üzvlər:}\\[3pt]
\begin{tabular}{@{}p{68mm}p{96mm}@{}}
  Kərimli Kənan Lətif oğlu & --- institutun direktor müavini \\
  Məmmədov Rəşad İlham oğlu & --- Elmi-pedaqoji tədqiqatlar mərkəzinin rəhbəri \\
  Quliyeva Nərminə Fikrət qızı & --- Qiymətləndirmə, təhlil və monitorinq mərkəzinin rəhbəri \\
  İsmayılov Tural Mübariz oğlu & --- Funksional şöbələrin müdiri \\
  Əliyeva Sevinc Kamran qızı & --- Funksional şöbələrin baş mütəxəssisi \\
  Rzayeva Günel Azər qızı & --- Funksional şöbələrin mütəxəssisi \\
  Babayev Elçin Sərvər oğlu & --- Təhsil texnologiyaları mərkəzinin aparıcı mütəxəssisi \\
  \adxetti & --- institutun baş mühasibi \\
  \adxetti & --- mühasib \\
\end{tabular}
\end{flushleft}

\vspace{20pt}

\begin{flushleft}
\begin{tabular}{@{}p{60mm}p{45mm}p{62mm}@{}}
  Direktor əvəzi & \imzaxetti & Əliyev Elnur Qəzənfər oğlu \\
\end{tabular}
\end{flushleft}

\end{document}
